\section{Заключение}
\label{sec:conclusion}

Настоящата дипломна работа представи проектирането, реализацията и анализа на библиотека за обектно-релационно съпоставяне по време на компилация за C++, ориентирана към работа както с релационни, така и с нерелационни бази данни. В центъра на разработката стои идеята, че структурата на заявката, логическият модел на данните и съществена част от ограниченията върху употребата на целевите системи могат да бъдат описани и валидирани чрез типовата система на езика.

С това работата адресира един конкретен проблем от съвременната практика: разминаването между силно типизиран приложен код и хетерогенните модели за съхранение, които този код трябва да използва. Класическите ОРС решения са най-силни в рамките на релационния свят, а подходите, ориентирани изцяло към драйвера (ориентирани към драйвера), често прехвърлят твърде голяма част от отговорността за коректността към времето за изпълнение. Настоящата разработка показа, че между тези два полюса съществува жизнеспособна архитектурна позиция --- ОРС слой по време на компилация, който е едновременно практичен, формално структуриран и достатъчно честен по отношение на различията между свързващите компоненти.

\subsection{Обобщение на решената задача}

Поставената задача не беше просто да се изгради още един удобен конструктор на заявки (заявка конструктор). Целта беше да се проектира библиотека, в която:

\begin{enumerate}[label=\arabic*.]
    \item логическият модел на данните се описва чрез C++ типове;
    \item заявките се представят като типизирано междинно представяне, а не като низове;
    \item целевите системи се интегрират чрез формален договор на конектора;
    \item различията между моделите за съхранение се управляват чрез механизми за възможности и ограничения;
    \item коректността се подкрепя едновременно от проверки по време на компилация и от тестова инфраструктура.
\end{enumerate}

В рамките на дипломната работа всяка от тези подцели беше адресирана с конкретни архитектурни решения и решения по реализацията. Слоят на същностите въведе типизирани свойства и отношения, политика за съхранение чрез \code{store\_as} и характеристики за именуване на логическите контейнери. Конструкторът на заявки беше организиран около типизирано междинно представяне. Конекторният слой беше формализиран чрез \code{connector\_trait<DB>}, а стратегията за верификация бе подкрепена от модулни и интеграционни тестове за множество свързващи компоненти.

\begin{figure}[H]
\centering
\begin{tikzpicture}[node distance=1.35cm and 1.15cm]
    \node[ormaccent] (model) {Логически модел\\типове същности + връзки};
    \node[ormblock, right=of model] (ir) {Типизирано междинно представяне\\предикати, проекции, заместители};
    \node[ormblock, below=of $(model)!0.5!(ir)$] (конектор) {Договор на конектора\\възможности + материализация};
    \node[ormblock, below left=of конектор] (tests) {Верификация\\модулни + интеграционни тестове};
    \node[ormblock, below right=of конектор] (свързващ компонентs) {Изпълнение върху\\множество целеви системи};

    \draw[ormline] (model) -- (ir);
    \draw[ormline] (model) -- (конектор);
    \draw[ormline] (ir) -- (конектор);
    \draw[ormline] (конектор) -- (tests);
    \draw[ormline] (конектор) -- (свързващ компонентs);
\end{tikzpicture}
\caption{Синтез на основната конструкция на разработената библиотека}
\label{fig:thesis_synthesis}
\end{figure}

Фигура~\ref{fig:thesis_synthesis} обобщава цялостната логика на дипломната работа. Тя показва, че библиотеката не е съвкупност от отделни техники, а последователна система, в която логическият модел, междинното представяне на заявките, договорът на конектора, верификацията и изпълнението върху множество компоненти се поддържат взаимно.

\subsection{Обобщение на постигнатите резултати}

Разработката показа, че C++ предоставя достатъчно силни механизми --- шаблони, \code{constexpr}, концепти (\code{concepts}) и специализации на характеристики --- за изграждане на ОРС слой, в който заявките не се описват като низове. Вместо това те се представят чрез типизирано междинно представяне, което може да бъде анализирано по време на компилация и материализирано по различен начин според целевата система.

На равнище модел на данните библиотеката демонстрира как \code{property}, \code{relationship}, \code{store\_as} и \code{table\_name\_trait} могат да играят ролята на единен логически слой. Този слой не е ограничен до релационната интерпретация, а служи като общо описание на данните и връзките между тях. Именно това позволява една и съща структура на приложно ниво да бъде пренасяна към таблици, документи, ключове и стойности, графови възли или широки колонни редове.

На равнище архитектура характеристиката на конектора (\code{connector\_trait}) е формализирана като централен механизъм за разширяемост. Именно чрез нея библиотеката адаптира едно и също междинно представяне към SQL, BSON-подобни филтри, Redis команди, Cypher и CQL. Механизмът на възможностите (възможност) допълва тази архитектура, като прави недопустимите операции видими още при компилация.

На равнище верификация разработката се опира на систематичен набор от модулни и интеграционни тестове, които покриват слоя на същностите, конструктора на заявки, механизма за подготвени заявки, миграциите, нишковата безопасност и конкретните конектори. По този начин архитектурните твърдения са подкрепени с реални доказателства от кода и тестовата инфраструктура.

\subsection{Научни и приложни приноси}

Първият принос на работата е демонстрацията, че практически полезна библиотека за обектно-релационно съпоставяне по време на компилация може да бъде реализирана в стандартен C++. Това се постига без външна кодогенерация, без виртуален интерфейс на конектора и без принуждаване на потребителя да напуска идиоматичния C++ модел.

Вторият принос е в изясняването на връзката между C++ описанието на същността и физическата организация на данните. Показано е, че C++ моделът може да бъде третиран като първичен носител на логическата схема, докато конкретната целева система определя начина на материализация. Тази постановка има стойност не само за ОРС библиотеките, а и по-общо за системите, които се опитват да комбинират силна статична типизация с хетерогенни източници на данни.

Третият принос е формализирането на договор на конектора, който обединява разширяемост и строга дисциплина по време на компилация. Този договор прави възможно надграждането на библиотеката с нови системи за съхранение, без да се разрушават вече съществуващите гаранции в програмния интерфейс за заявки. Точно тук работата се различава от подходи, които предлагат общ интерфейс, но скриват критичните разлики между целевите системи.

Четвъртият принос е показването, че верификацията на подобна библиотека трябва да бъде многостепенна. Не е достатъчно само ``да има тестове'' или само ``да има проверки по време на компилация''. Необходима е комбинация от двете: ограничения на ниво концепти, утвърждаване на възможности, тестове за рендиране на заявки, интеграционни сценарии на живо и проверки на миграциите и нишковата безопасност.

\begin{table}[H]
\centering
\small
\caption{Връзка между целите на дипломната работа и реализираните резултати}
\label{tab:goal_result_mapping}
\begin{tabularx}{\textwidth}{L{3.1cm}L{3.2cm}Y}
\toprule
\textbf{Цел} & \textbf{Реализиран механизъм} & \textbf{Доказателствена опора} \\
\midrule
Типизиран модел на данните & \code{property}, \code{relationship}, \code{table\_name\_trait} & заглавни файлове на същности и модулни тестове за свойства/връзки \\
Типизиран слой за заявки & конструктори на заявки и типизирано междинно представяне & заглавни файлове на заявки, \path{tests/модулни/test_заявка.cpp}, \path{tests/интеграционни/test_mockdb.cpp} \\
Разширяем слой за множество целеви системи & договор на конектора и labelи за възможности & заглавни файлове на конектори и платформено-специфични модулни тестове / тестове на живо \\
Посока, съобразена със схемата & \code{migrate<DB>} и DDL куки & \path{tests/модулни/test_schema_migration.cpp} \\
Безопасна експлоатационна употреба & помощници за нишкова безопасност и транзакции & \path{tests/модулни/test_нишка_safety.cpp} \\
\bottomrule
\end{tabularx}
\end{table}

Таблица~\ref{tab:goal_result_mapping} показва, че дипломната работа не остава на равнището на декларативен архитектурен документ (архитектурен документ). За всяка основна цел има конкретен реализиран механизъм и поне един пряк технически корелат в хранилището.

\subsection{Основни ограничения и граници на обобщението}

Работата също така ясно показва, че ОРС абстракцията за множество целеви системи има граници. Не всеки модел за съхранение предлага едни и същи операции, а следователно общото междинно представяне не може да бъде интерпретирано еднакво навсякъде. Именно затова механизмите за възможности и ограничения са толкова важни. Те не отслабват библиотеката, а я правят коректна спрямо различните системи.

Допълнително, макар архитектурната основа за миграции, нишкова безопасност и по-ниско ниво на изпълнение вече да съществува, тези подсистеми остават естествено поле за по-нататъшно развитие и емпирично разширяване. Това означава, че работата трябва да бъде оценявана като завършена архитектурна демонстрация и демонстрация на реализацията, но не и като последна дума по темата за достъп до данни по време на компилация.

От академична гледна точка е важно да се подчертае и още една граница. Настоящата разработка не доказва, че ОРС по време на компилация е универсално най-добрият подход за всички приложения. Тя доказва нещо по-точно: че за определен клас системи --- системи с нужда от ясна статична структура, висока типова безопасност (типова безопасност) и контролирана разширяемост към множество целеви системи --- този подход е жизнеспособен и инженерно защитим.

\subsection{Значение за практиката и за бъдещите изследвания}

Практическата стойност на разработката се състои в това, че библиотеката вече осигурява стабилна основа за по-нататъшно инженерно развитие. Описанието на същностите, междинното представяне на заявките, договорът на конектора и стратегията за верификация са достатъчно ясно формулирани, за да позволят дисциплинирано добавяне на нови целеви системи и експлоатационни възможности.

Изследователската стойност е свързана с по-общия въпрос за ролята на типовата система в моделирането на достъпа до данни. Работата показва, че типовата система на C++ може да служи не само за описание на данни и алгоритми, но и за формализиране на архитектурни граници, допустимост на операции и преход между логически и физически слоеве. Това поставя библиотеката в по-широк контекст, който надхвърля ОРС темата в тесния смисъл.

\subsection{Заключителни бележки}

Направеният анализ потвърждава, че типовата система на C++ е достатъчно мощна, за да служи не само за описване на данни и алгоритми, но и за формализиране на достъпа до персистентни системи. Комбинацията от статичен модел на същностите, междинно представяне, конекторен слой, базиран на характеристики, и систематична верификация показва устойчив път към библиотеки, които са едновременно изразителни, безопасни и надграждаеми.

По този начин дипломната работа постига както инженерна, така и изследователска цел: изгражда работеща библиотека и едновременно формулира принципи, които могат да бъдат използвани при бъдещи разработки на слоеве за достъп до данни по време на компилация. Най-същественият извод е, че подходът за обектно-релационно съпоставяне по време на компилация не е просто теоретично упражнение, а реална инженерна алтернатива за системи, които изискват строгост, яснота и разширяемост при работа с множество модели за съхранение.
